\section{Limitations and predictions}
\label{sec:limits}

\subsection{Limitations}

\paragraph{Three cases, chosen by the author.} The cases were selected because the
vocabulary fits them. They are evidence that the vocabulary of LEM can be carried to
these three settings, and no more. Cases in which typing a crossing adds nothing, or in
which the kinds proposed here do not apply, were not searched for.

\paragraph{A parallel, not an identity.} The three roles of a map in
Section~\ref{sec:three-kinds} and the six edge types of Table~\ref{tab:edges} overlap but
were arrived at separately. The paper does not give one classification that covers both,
and does not claim that either list is complete or that its members exclude one
another.

\paragraph{The mechanisation is partial.} The Lean package states the typed crossing
and works it out for two maps from the first case. The refinement case is not
formalised, the reasoning-system model is not expressed as a crossing, and the package
is not connected to Mathlib (Section~\ref{sec:mech-limits}). That the kind of a
crossing is not determined by its relation holds in the formal model by construction,
and so adds no evidence for the claim.

\paragraph{The kind is declared, not decided.} The rule asks that the kind of a crossing
be declared. It does not supply a procedure that decides it. LEM proposes tests for edge
types, two of which are used in Section~\ref{sec:duration}; they are not formalised, and
their agreement between independent raters has not been measured. Where two modellers
would declare different kinds for one crossing, the rule exposes the disagreement and
does not resolve it.

\paragraph{The results are elementary and rest on a model.} The weight of
Section~\ref{sec:kernel} is carried by Definition~\ref{def:rs}, not by the proofs. The
model treats faithfulness as a binary relation between a trace and a mechanism,
presupposes that mechanisms can be individuated, and takes production to be a function.
Each is a simplification (Remark~\ref{rem:scope}), and a graded notion of faithfulness
is not covered. The deontic model of Section~\ref{sec:deontic} is smaller still: it has
no deadlines and no conditions of applicability.

\paragraph{No empirical evaluation.} The paper reports no experiments. In particular it
does not show that real models contain two computations that write the same trace with
different verdicts, nor how often, which is what determines how much
Corollary~\ref{cor:split} matters in practice.

\paragraph{The source of the vocabulary.} LEM is a method proposal by the same author,
developed for risk modelling and published as a preprint. Its own validation is open,
and this paper inherits that status for everything it takes from LEM.

\paragraph{An outside view.} The paper addresses three research communities, and the
author's own field is risk modelling. The accounts given here of Mathlib, of
verification practice and of interpretability research rest on the literature cited, not
on work inside those communities. Corrections from them are invited.

\subsection{Two predictions}

The argument yields two predictions that existing benchmarks of automated theorem
proving could test. They are independent, and the second is the one that bears on the
rule.

\begin{prediction}[Crossings are where provers fail]
Among the failures of automated provers on formal proof tasks, proof steps that cross a
boundary (a cast, a change of representation, a step from an operational semantics to an
invariant) are over-represented, relative to their share of all steps and compared with
steps inside one structure.
\end{prediction}

\begin{prediction}[Declaring the kind helps beyond supplying the lemmas]
Give a prover the preservation lemmas that are relevant to the crossings in a task. A
prover that is given, in addition, a declaration of the kind of each crossing fails
less often at crossings than one that is given the same lemmas without it.
\end{prediction}

No measurement is cited or reported here. The first prediction requires a classification
of proof steps into crossing and non-crossing ones, fixed before the failure data are
examined, and would be refuted by a failure rate at crossings that does not differ from
the rate elsewhere. The second needs the baseline it names. Without it, an improvement
would show only that supplying useful lemmas helps a prover, which is not in doubt. If
the declaration of kind adds nothing over the lemmas, the first component of
Definition~\ref{def:typed} has no measurable use in this setting, and the paper's case
for it would rest on the diagnostic uses of Sections~\ref{sec:refinement}
and~\ref{sec:reasoning} alone.
